(1) 设TM模式磁场为 \(H_y^{(d)}=H_0e^{-\kappa_d z}e^{i\beta x-i\omega t}\)（\(z>0\)），\(H_y^{(m)}=H_0e^{+\kappa_m z}e^{i\beta x-i\omega t}\)（\(z<0\)）。各侧满足 \(\kappa_j^2=\beta^2-\varepsilon_jk_0^2\)，并选 \(\operatorname{Re}\kappa_j>0\) 使能量局域。由 \(\nabla\times\mathbf H=-i\omega\varepsilon_0\varepsilon_j\mathbf E\) 得 \(E_x^{(d)}=-i\kappa_dH_y^{(d)}/(\omega\varepsilon_0\varepsilon_d)\)、\(E_x^{(m)}=+i\kappa_mH_y^{(m)}/(\omega\varepsilon_0\varepsilon_m)\)；磁场和电场均有横纵分量，不是TE表面波。

(2) 无自由面电流时 \(H_y\) 连续，切向 \(E_x\) 连续给 \(\kappa_d/\varepsilon_d+\kappa_m/\varepsilon_m=0\)。结合 \(\kappa_j^2=\beta^2-\varepsilon_jk_0^2\)，平方并排除不满足衰减号的根，得 \(\boxed{\beta=k_0\sqrt{\varepsilon_d\varepsilon_m/(\varepsilon_d+\varepsilon_m)}}\)。对 \(\varepsilon_d>0\) 的无损金属界面，实且高于介质光线 \(\beta>k_0\sqrt{\varepsilon_d}\) 的束缚TM模要求 \(\varepsilon_m<-\varepsilon_d\)。TE波的边界式不能在两侧普通 \(\mu>0\) 时产生相同束缚解。

(3) 直接代入 \(\beta\) 可得 \(\kappa_d=k_0\varepsilon_d/\sqrt{-(\varepsilon_m+\varepsilon_d)}\)、\(\kappa_m=k_0(-\varepsilon_m)/\sqrt{-(\varepsilon_m+\varepsilon_d)}\)（取正实根），穿透长度即各自倒数。对 \(\varepsilon_d=1,\varepsilon_m=1-\omega_p^2/\omega^2\)，当 \(\beta\to\infty\) 必有 \(\varepsilon_m+\varepsilon_d\to0^-\)，故 \(\boxed{\omega_{\rm sp}=\omega_p/\sqrt2}\)。在 \(\omega\ll\omega_p\) 时 \(\varepsilon_m\simeq-\omega_p^2/\omega^2\)，展开 \(\beta\simeq(\omega/c)[1+\omega^2/(2\omega_p^2)+O(\omega^4/\omega_p^4)]\)，渐近真空光线。

(4) 在无损Drude模型中 \(\beta\) 随 \(\omega\to\omega_{\rm sp}^-\) 发散，因此 \(v_g=(d\beta/d\omega)^{-1}\to0\)，表面模强烈局域。实际金属有 \(\operatorname{Im}\varepsilon_m>0\)，于是 \(\beta\) 变为复数，沿传播方向出现有限衰减，严格的无穷大 \(\beta\) 不可实现。无损色散介质的电能密度应采用Brillouin形式 \(\frac14\varepsilon_0\partial_\omega[\omega\varepsilon(\omega)]|E|^2\)，Drude中 \(\partial_\omega(\omega\varepsilon_m)=1+\omega_p^2/\omega^2>0\)，避免把负介电常数误判为负的物理储能。
